\section*{Chapter \ref{ch:iv:offers-compensation-and-non-competes} --- Offers, Compensation and Non-Competes}

\begin{solution}{iq:iv:offers-compensation-and-non-competes:1}
In order: base salary and currency; whether the first bonus is guaranteed and how later bonuses are decided;
the deferral share, instrument, vesting schedule and leaver rules (including malus and clawback); one-off
payments (sign-on, buy-out) and their repayment conditions; start date and notice period; restrictive
covenants (length, scope, whether paid, whether the firm can waive them); and the conditions precedent. Without
the leaver rules and the covenants, two offers with the same headline cannot be compared.
\iqlookfor{the full set of terms that change an offer's value, and why the leaving terms matter.}
\end{solution}

\begin{solution}{iq:iv:offers-compensation-and-non-competes:2}
Fifteen of the 24 months remain: $60\,000 \times 15/24 = 37\,500$.
\iqlookfor{pro-rata arithmetic, and reading the clause's clock correctly.}
\end{solution}

\begin{solution}{iq:iv:offers-compensation-and-non-competes:3}
Cash at the bonus date $0.6 \times 200\,000 = 120\,000$; deferred $80\,000$ in three instalments of about
$26\,667$ at 12, 24 and 36 months. At 18 months one instalment has been paid and two are unvested:
$53\,333$ is forfeited, unless the new employer buys it out (One Quant Book~16, chapter~10, on deferral
buy-outs), which is the item to negotiate.
\iqlookfor{the deferral schedule, what is at risk on leaving, and the buy-out as the remedy.}
\end{solution}

\begin{solution}{iq:iv:offers-compensation-and-non-competes:4}
A is worth $1\,070 - 400\pi$ and B $960 - 60\pi$. At $\pi = 0.2$: A 990, B 948, so A. At $\pi = 0.5$: A
870, B 930, so B. Indifference at $\pi = 110/340 = 11/34 \approx 0.32$. The honest answer adds that $\pi$ is
not known, and that offer A's value is far more sensitive to it (400 against 60 a unit of probability): A is
the bet that you will stay.
\iqlookfor{expected value linear in the leaving probability, the break-even, and sensitivity.}
\end{solution}

\begin{solution}{iq:iv:offers-compensation-and-non-competes:5}
Thank the firm, say that you are very interested and that you are completing one other process whose \omterm{def:iv:the-final-round-and-trading-games:final}{final
round} is next week, and ask for an answer date after it (give the date). Offer something in return: a call on
a named day, an early start. If the firm refuses, value the offer against your estimate of the other process
(\cref{ch:iv:the-process-by-firm-type}, the accept-or-wait method) and decide; do not accept with the intention
of reneging.
\iqlookfor{a specific request for time, honesty about the other process, and a decision rule if refused.}
\end{solution}

\begin{solution}{iq:iv:offers-compensation-and-non-competes:6}
Paid garden leave of six months costs the difference between the next job's pay and the garden-leave pay for
six months: $0.5 \times (400 - 250) = 75$. The unpaid twelve-month non-compete costs a full year of the next
job's pay: 400. Negotiate the non-compete: ask for it to be shortened, paid, or limited in scope to the
competitors that matter, or ask for a sign-on that compensates it.
\iqlookfor{pricing covenants as lost pay, and knowing that paid and unpaid restrictions differ.}
\end{solution}

\begin{solution}{iq:iv:offers-compensation-and-non-competes:7}
On the number: give a range anchored on evidence, not on hope (``for this role in this city, total compensation
of X to Y, based on the published range and my other conversations''), or ask for the firm's range first where
it must publish one. On current pay: in some places the firm may not ask
(\cref{dat:iv:offers-compensation-and-non-competes:law}); elsewhere, you may decline politely and redirect to
what the role is worth, or answer truthfully including the deferrals you would forfeit, which the firm will
have to buy out. Never inflate it: it can be verified.
\iqlookfor{anchoring on market evidence, awareness of the law, and truthfulness.}
\end{solution}

\begin{solution}{iq:iv:offers-compensation-and-non-competes:8}
Compare the two offers as cash flows over the same horizon (base, expected bonus, deferrals and their
forfeiture, covenants), with the same probability-of-leaving analysis; check what is written (a 20\% base rise
is written, a larger bonus usually is not). The quantity you cannot compute is whether the reasons you started
looking are removed: if they were not about money, the \omterm{def:iv:offers-compensation-and-non-competes:exploding}{counteroffer} buys time, and the next move starts from a
position in which your employer knows you were leaving. Decide on the written terms and on those reasons, and
honour the commitment you have made to the new firm unless you have not yet accepted.
\iqlookfor{a like-for-like valuation, the distinction between written and promised pay, and the non-financial
question.}
\end{solution}

\begin{solution}{iq:iv:offers-compensation-and-non-competes:9}
Asking $a$ is accepted with probability $(350 - a)/100$, so the expected value is $260 + (a - 260)(350 -
a)/100$, a parabola maximised at $a = (260 + 350)/2 = 305$, with acceptance probability 0.45 and expected
value $260 + 45^2/100 = 280.25$. The model leaves out that negotiation is not one shot (a counterproposal is the
usual reply), that asking may change the firm's view of you, and that the firm's walk-away value is not
uniform; its useful lesson is that the best ask is well above your alternative and well below the firm's
maximum.
\iqlookfor{an optimisation over the ask with an acceptance probability, and the model's limits.}
\end{solution}

\begin{solution}{iq:iv:offers-compensation-and-non-competes:10}
It forbids working in any competing business, in any role, anywhere the group operates, for a year, unpaid
unless another clause pays it. Ask: whether it is paid, whether the firm will confirm in writing which
competitors it has in mind, whether it can be waived and in practice is, how it interacts with garden leave
and notice (whether the periods run concurrently), and which law governs the contract. Enforceability depends
on the law of the place: in California such a clause is void and may not be enforced
(\cref{dat:iv:offers-compensation-and-non-competes:noncompete}); in the United Kingdom the starting point is that a clause in restraint
of trade is unenforceable unless the employer shows that it is reasonable, which a court decides if the clause
is challenged (as the government's 2025 working paper puts it), and a twelve-month worldwide ban on any
competing business is hard to show reasonable for most roles. Take legal advice before signing.
\iqlookfor{reading the scope of a clause, the right questions, and the role of the governing law.}
\end{solution}
